\documentstyle[% 


righttag% 


,11pt% 


,amssymb% 


,russian%               remove in Eng. 


,izvru%                 Set izven% in Eng. 


]{amsart} 


 


%       Math definitions 


 


\newcommand{\BH}{\text{\it Б}^{\text{\it г}}_4} 


% NB change in eng.! ^^^^^^^^^^^^^^^ 


\newcommand{\tts}[1]{} 


 


%\hoffset=-15mm%                  For Eng. 


  


\setcounter{page}{83} 


\god{1999} 


\nomer{9~(448)} %                        Remove in Eng. 


%\nomer{Vol.\,43, No.\,9}%               Set in Eng. 


%\str{\pageref{begin}--\pageref{end}}%  Set in Eng. 


\udk{514.754} 


 


\author{\it Л.Н.\,Ромакина} 


% NB:   ^^ remove in Eng. 


\title{Теория кривых в биаффинно-флаговом  


пространстве гиперболического типа} 


 


\date{} 


% \pagestyle{myheadings}%         Set in Eng. 


\pagestyle{plain} %              Remove in Eng. 


\begin{document} 


% \thispagestyle{plain}%          Set in Eng. 


\maketitle 


\label{begin} 


 


Биаффинно-флаговым пространством гиперболического  


типа $\BH$ назовем четырехмерное аффинное пространство, фундаментальной 


группой которого является группа $G$ аффинных  


преобразований, оставляющая в несобственной гиперплоскости  


инвариантным абсолют, состоящий из абсолюта биаксиального  


пространства гиперболического типа [1] с выделенной прямой  


абсолютной конгруэнции. Пространство $\BH$ является вещественной 


реализацией двойной бифлаговой плоскости.


Его группа движений зависит от девяти параметров. В данной  


работе строится теория кривых пространства $\BH$. 


\medskip 


 


{\bf 1.} При подходящем задании абсолюта группа $G$ пространства 


$\BH$ изоморфна группе матриц вида  


$$ 


\begin{pmatrix} 


a_1 & b_1 & a_2 & b_2 & h_1 \\ 


\varepsilon b_1 & \varepsilon a_1 & \varepsilon b_2 & 


\varepsilon a_2 & h_2 \\ 


0 & 0 & a_3 & b_3 & h_3 \\ 


0 & 0 & \varepsilon b_3 & \varepsilon a_3 & h_4 


\end{pmatrix} 


, 


$$ 


где $\varepsilon=\pm 1$, $(a_3)^2-(b_3)^2=1$. 


 


Относительно группы $G$ в $\BH$ инвариантно скалярное произведение 


векторов $\ov a(a_i)$, $\ov b(b_i)$ ($i=1,2,3,4$), которое определяется 


вырожденной квадратичной 


формой ранга 2 индекса 1, 


вычисляется по формуле 


$$ 


\ov a\ov b=a_3b_3-a_4b_4, 


$$ 


а также гиперболическое измерение углов между 


$m$-плоскостями ($m=1,2,3$). Таким образом, $\BH$ 


является пространством 


с обобщенной проективной метрикой [2]. 


 


Два вектора назовем ортогональными,если их скалярное  


произведение равно нулю. Если направления ортогональных  


векторов определяют в несобственной гиперплоскости точки  


прямой абсолютной конгруэнции, то назовем их сильно ортогональными. 


Векторы, определяющие точки особой прямой абсолюта $P_1$, назовем 


изотропными. 


 


Очевидно, изотропный вектор ортогонален любому неизотропному 


вектору. Для двух изотропных векторов $\ov a(a_i)$, $\ov b(b_i)$ 


определим скалярное произведение второго рода, которое инвариантно 


относительно группы $G$ и вычисляется по формуле 


$$ 


(\ov a\ov b)=a_1b_1-a_2b_2. 


$$ 


Если скалярное произведение второго рода двух изотропных  


векторов равно нулю, то назовем эти векторы изотропными ортогональными. 


 


Будем рассматривать координатные реперы $R\{O,\ov e_i\}$,


обладающие следующими инвариантными 


относительно преобразований группы $G$ свойствами: 


\begin{itemize} 


\item[1)] $O$ --- произвольная точка пространства; 


\item[2)] векторы $\ov e_1$, $\ov e_2$ изотропные ортогональные, т.\,е. 


$\ov e^2_1=\ov e^2_2=0$ и $(\ov e_1\,\ov e_2)=0$; 


\item[3)] векторы $\ov e_3$, $\ov e_4$ сильно ортогональные, причем 


один из  


них единичный, другой мнимоединичный, т.\,е. 


$\ov e_3\,\ov e_4=0$, $(\ov e_3\,\ov e_4)=0$, 


$\ov e^2_3=-\ov e^2_4=\pm 1$. 


\end{itemize} 


 


С каждой точкой пространства свяжем семейство канонических 


реперов, обладающих инвариантными свойствами 1)--3).  


 


\begin{thm}


Канонический репер $R\{M,\ov n_i\}$ пространства $\BH$ однозначно 


определяется заданием точки $M$ и направлений векторов 


$\ov n_3$ и $\ov n=\ov n_1+\ov n_4$. 


\end{thm} 


 


{\bf 2.} Построим канонический репер $R\{M,\ov n_i\}$ в произвольной  


точке $M$ пространства $\BH$. Его инфинитезимальное смещение  


определяется уравнениями [3] 


$$ 


d\ov M=\omega^i\ov n_i,\quad 


d\ov n_i=\omega^j_i\ov n_j,\quad i,j=1,2,3,4. 


$$ 


 


Уравнения структуры группы движений пространства $\BH$ 


получим, учитывая


$$ 


D\omega^i=[\omega^k\omega^i_k],\quad 


D\omega^j_i=[\omega^k_i\omega^j_k],\quad i,j,k=1,2,3,4. 


$$ 


Формы $\omega^j_i$ предыдущих уравнений связаны условиями 


\begin{gather*} 


\omega^3_1=\omega^3_2=\omega^4_1=\omega^4_2= 


\omega^3_3=\omega^4_4=\omega^4_3-\omega^3_4=0;\\ 


\omega^1_2-\omega^2_1=\omega^1_1-\omega^2_2= 


\omega^1_3-\omega^2_4=\omega^2_3-\omega^1_4=0. 


\end{gather*} 


 


{\bf 3.} Зададим кривую в пространстве $\BH$ уравнением 


\begin{equation} 


\ov r=\ov r(t), 


\end{equation} 


где $\ov r(t)$ имеет непрерывные производные до третьего порядка  


включительно. 


Назовем точку кривой (1) особой, если выполняется одно  


из следующих условий: 


\begin{itemize} 


\item[1)] производные до третьего порядка ее радиуса-вектора  


линейно зависимы; 


\item[2)] ее соприкасающаяся 2-плоскость пересекает абсолютную прямую; 


\item[3)] ее соприкасающаяся 3-плоскость содержит абсолютную 


прямую. 


\end{itemize} 


 


В каждой неособой точке кривой (1) можно построить репер 


третьего порядка, но он не будет каноническим репером  


пространства $\BH$. Поэтому в точке $M$ кривой этого пространства 


строим два репера, один канонический $R\{M,\ov n{}^{\,i}\}$, другой --- репер 


Френе $R\{M,\ov t{}^{\,i}\}$. Связь между ними характеризует  


 


\begin{thm} В каждой неособой точке кривой $(1)$ в пространстве $\BH$ 


существует единственный канонический репер и репер  


Френе, который выражается через канонический с помощью  


одного параметра. 


\end{thm} 


 


\begin{thm} В пространстве $\BH$ имеют место формулы типа  


формул Френе, содержащие три независимых параметра. 


\end{thm} 


 


\begin{thebibliography}{99} 


 


\bibitem{1} 


Широков А.П. {\em Геометрия обобщенных биаксиальных пространств} // 


Учен. зап. Казанск. ун-та. -- 1954. -- Т.\,114. -- Вып.\,2. -- 


С.\,123--166. 


\bibitem{2} 


Киотина Г.В. {\em Пространства с обобщенной проективной метрикой}. 


Учеб. пособие. -- Рязань: Изд-во Рязанск. гос. пед. ун-та, 


1981. -- 102 с. 


\bibitem{3} 


Фиников С.П. {\em Метод внешних форм Картана}. -- М.--Л.: Гостехиздат, 


1948. -- 432 с. 


 


\end{thebibliography} 


 


\vskip 2cm 


 


\postup{\it Саратовский государственный}{$\qquad$\it Поступили} 


\postup{\it университет}{{\it полный текст} 24.02.1998} 


\postup{}{{\it краткое сообщение} 23.09.1998} 


\label{end} 


\end{document} 


